\documentclass[11pt,a4paper]{article}

\usepackage[utf8]{inputenc}
\usepackage[T1]{fontenc}
\usepackage[spanish,es-noshorthands]{babel}
\usepackage[margin=2.4cm]{geometry}
\usepackage{amsmath,amssymb}
\usepackage{xcolor}
\usepackage{tikz}
\usetikzlibrary{arrows.meta,positioning,shapes.geometric,fit,backgrounds,calc}
\usepackage{booktabs}
\usepackage{tcolorbox}
\usepackage{listings}
\usepackage{enumitem}
\usepackage{fancyhdr}
\usepackage{titlesec}
\usepackage[hidelinks]{hyperref}

% ---- Colors ----
\definecolor{ghblue}{HTML}{1F4E79}
\definecolor{ghteal}{HTML}{2C7A7B}
\definecolor{ghgray}{HTML}{4A5568}
\definecolor{ghlight}{HTML}{EBF4FF}
\definecolor{ghgreen}{HTML}{22803B}
\definecolor{ghred}{HTML}{9B2C2C}
\definecolor{ghamber}{HTML}{B7791F}
\definecolor{codebg}{HTML}{F5F7FA}

% ---- Section styling ----
\titleformat{\section}{\color{ghblue}\Large\bfseries}{\thesection.}{0.6em}{}
\titleformat{\subsection}{\color{ghteal}\large\bfseries}{\thesubsection}{0.6em}{}

% ---- Header/footer ----
\pagestyle{fancy}
\fancyhf{}
\fancyhead[L]{\small\color{ghgray}SOT en GraphHunter}
\fancyhead[R]{\small\color{ghgray}Teoría del Observador Suficiente}
\fancyfoot[C]{\small\color{ghgray}\thepage}
\renewcommand{\headrulewidth}{0.4pt}

% ---- Listings ----
\lstset{
  basicstyle=\ttfamily\footnotesize,
  backgroundcolor=\color{codebg},
  frame=single,
  rulecolor=\color{ghgray!40},
  breaklines=true,
  columns=fullflexible,
  keepspaces=true,
  showstringspaces=false,
  xleftmargin=6pt,xrightmargin=6pt,
  aboveskip=8pt,belowskip=8pt,
}

% ---- Boxes ----
\newtcolorbox{ideabox}[1][]{colback=ghlight,colframe=ghblue,
  fonttitle=\bfseries,title={#1},boxrule=0.8pt,arc=2pt,left=6pt,right=6pt,top=4pt,bottom=4pt}
\newtcolorbox{keybox}[1][]{colback=ghgreen!8,colframe=ghgreen,
  fonttitle=\bfseries,title={#1},boxrule=0.8pt,arc=2pt,left=6pt,right=6pt,top=4pt,bottom=4pt}
\newtcolorbox{warnbox}[1][]{colback=ghamber!10,colframe=ghamber,
  fonttitle=\bfseries,title={#1},boxrule=0.8pt,arc=2pt,left=6pt,right=6pt,top=4pt,bottom=4pt}

\newcommand{\term}[1]{\textcolor{ghblue}{\textbf{#1}}}
\newcommand{\code}[1]{\texttt{\small #1}}

% ---- Title ----
\title{\vspace{-1.5cm}
  {\color{ghblue}\Huge\bfseries Teoría del Observador Suficiente\\[2mm]
  \Large en GraphHunter}\\[4mm]
  {\color{ghgray}\large Cómo un cazador de amenazas sabe cuándo \emph{puede confiar} en lo que ve\\
  \normalsize Guía didáctica, paso a paso}}
\author{\color{ghgray}Documentación de arquitectura --- GraphHunter}
\date{\color{ghgray}\today}

\begin{document}
\maketitle
\thispagestyle{fancy}

\begin{abstract}
\noindent Este documento explica, de forma didáctica y paso a paso, cómo la
\emph{Teoría del Observador Suficiente} (SOT, por \emph{Sufficient Observer
Theory}) organiza la ingesta y la caza de amenazas en GraphHunter. La idea
central es simple pero poderosa: antes de responder ``¿hay un ataque?'', un
observador honesto debe poder responder ``¿tengo suficiente información para
que mi respuesta sea confiable, o hay algo que no puedo ver?''. Recorreremos
qué es un observador, qué significa que sea \emph{suficiente}, qué es un
\emph{testigo de ceguera}, cómo está construida la tubería de datos de
GraphHunter, y seguiremos un evento real desde la fuente hasta la respuesta de
una hipótesis de caza.
\end{abstract}

\tableofcontents
\vspace{4mm}

% ======================================================================
\section{El problema: ver no es lo mismo que saber}

Imaginá un analista de seguridad frente a una consola. Escribe una hipótesis
de caza --- por ejemplo, \emph{``¿algún usuario se autenticó en un host y luego
ese host se comunicó con una IP externa sospechosa?''} --- y el sistema
responde: \textbf{``0 resultados''}.

\begin{ideabox}[La pregunta incómoda]
¿``0 resultados'' significa \emph{``no pasó''}, o significa \emph{``pasó, pero
no lo vi''}? Estas dos situaciones son operacionalmente idénticas para un
sistema ingenuo --- y peligrosamente distintas para el analista.
\end{ideabox}

Un firewall pudo haberse caído durante la ventana. Un log pudo llegar tarde y
quedar fuera del cursor de ingesta. Una tabla de Sentinel pudo dar timeout.
En todos esos casos, un sistema que solo dice ``0 resultados'' está
\textbf{mintiendo por omisión}: presenta datos incompletos como si fueran
exhaustivos.

SOT ataca exactamente este problema. No busca que el sistema \emph{nunca} sea
ciego --- eso es imposible en el mundo real --- sino que el sistema
\textbf{sepa cuándo es ciego} y lo \textbf{declare honestamente} en su
respuesta.

% ======================================================================
\section{Vocabulario SOT en cinco ideas}

SOT modela un sistema de observación con cinco elementos. Los presentamos con
su nombre formal y su traducción al mundo de GraphHunter.

\begin{center}
\renewcommand{\arraystretch}{1.35}
\begin{tabular}{@{}p{1.4cm}p{3.6cm}p{9.4cm}@{}}
\toprule
\textbf{Símbolo} & \textbf{Nombre} & \textbf{En GraphHunter} \\
\midrule
$X$ & Mundo real (estados) & Lo que \emph{realmente} pasó en la red: logins, conexiones, procesos. \\
$S$ & Señal / observación & Lo que las fuentes (Sentinel, Sysmon, FortiGate\ldots) \emph{reportan}. \\
$B$ & Familia de comportamientos & El catálogo de hipótesis de caza (30 patrones ATT\&CK) que queremos poder distinguir. \\
$C$ & Contextos & Condiciones bajo las que observamos: llegada tardía, caída de fuente, edición de mapping\ldots \\
$R$ & Respuesta del observador & La respuesta del \code{run\_hunt}: los caminos encontrados \emph{más} los metadatos de cobertura. \\
\bottomrule
\end{tabular}
\end{center}

\begin{keybox}[Definición de suficiencia]
Un observador es \term{suficiente} para una familia $B$ de comportamientos si,
a partir de lo que observa ($S$), puede \textbf{distinguir} cuál
comportamiento de $B$ ocurrió --- \emph{en todos los contextos $C$ declarados}.
Si existe algún contexto donde dos comportamientos distintos producen la
\emph{misma observación}, el observador es \textbf{insuficiente}, y esa
coincidencia se llama un \textbf{testigo de ceguera}.
\end{keybox}

\subsection{El testigo de ceguera, el corazón de SOT}

Un \term{testigo} (witness) es una prueba concreta de que el observador no
puede distinguir dos historias que \emph{deberían} distinguirse. Se describe
con tres piezas:

\begin{center}
\begin{tikzpicture}[node distance=3mm,every node/.style={font=\small}]
  \node[draw=ghblue,fill=ghlight,rounded corners,text width=4.3cm,align=center,minimum height=1.5cm] (l) {\textbf{Historia izquierda}\\ el evento llega \emph{tarde}};
  \node[draw=ghblue,fill=ghlight,rounded corners,text width=4.3cm,align=center,minimum height=1.5cm,right=1.6cm of l] (r) {\textbf{Historia derecha}\\ el evento llega \emph{puntual}};
  \node[draw=ghred,fill=ghred!8,rounded corners,text width=4.6cm,align=center,minimum height=1.5cm,below=1.1cm of $(l)!0.5!(r)$] (o) {\textbf{\code{same\_observed}}\\ la arista exportada es \emph{idéntica}};
  \draw[-{Latex[length=2mm]},ghgray] (l.south) -- (o.north west);
  \draw[-{Latex[length=2mm]},ghgray] (r.south) -- (o.north east);
  \node[right=0.5cm of o,text width=4.6cm,align=left,ghred] {\footnotesize $\Rightarrow$ el observador \textbf{no} puede razonar sobre llegada tardía: \term{testigo}.};
\end{tikzpicture}
\end{center}

Cerrar un testigo significa hacer que \code{same\_observed} \textbf{deje de ser
verdad}: cambiar el sistema para que las dos historias produzcan respuestas
distinguibles. Como veremos, GraphHunter cierra cada testigo con una pieza
concreta de código y lo \emph{prueba} con un test.

% ======================================================================
\section{La arquitectura: una tubería que preserva la evidencia}

El principio de diseño de GraphHunter es \textbf{``KQL no computa el grafo''}:
las fuentes (como Microsoft Sentinel) solo entregan \emph{eventos crudos}; toda
la correlación y la caza ocurren localmente, en un grafo de procedencia. Esto
desacopla al observador de la fuente y le da control total sobre la
suficiencia.

\begin{center}
\begin{tikzpicture}[
    node distance=6mm and 9mm,
    box/.style={draw=ghblue,fill=ghlight,rounded corners,text width=2.35cm,align=center,minimum height=1.15cm,font=\footnotesize\bfseries},
    every path/.style={-{Latex[length=2.4mm]},ghteal,thick}
  ]
  \node[box] (src) {Source\\Connector};
  \node[box,right=of src] (jr) {Raw Event\\Journal};
  \node[box,right=of jr] (nm) {Normalizer};
  \node[box,right=of nm] (gd) {Graph\\Delta};
  \node[box,below=of gd] (rt) {Incremental\\Graph Runtime};
  \node[box,left=of rt] (api) {Domain\\API};
  \node[box,left=of api] (mcp) {MCP\\/\,GUI};

  \draw (src) -- (jr);
  \draw (jr) -- (nm);
  \draw (nm) -- (gd);
  \draw (gd) -- (rt);
  \draw (rt) -- (api);
  \draw (api) -- (mcp);

  \node[below=1mm of src,text width=2.35cm,align=center,font=\scriptsize,ghgray] {KQL solo para \emph{traer} filas};
  \node[below=1mm of jr,text width=2.35cm,align=center,font=\scriptsize,ghgray] {append-only, rotado};
  \node[below=1mm of nm,text width=2.35cm,align=center,font=\scriptsize,ghgray] {sobre canónico};
  \node[above=1mm of gd,text width=2.35cm,align=center,font=\scriptsize,ghgray] {idempotente};
  \node[right=1mm of rt,text width=2.0cm,align=left,font=\scriptsize,ghgray] {el grafo local};
  \node[below=1mm of api,text width=2.35cm,align=center,font=\scriptsize,ghgray] {\code{run\_hunt} vive acá};
  \node[below=1mm of mcp,text width=2.35cm,align=center,font=\scriptsize,ghgray] {caza 100\% local};
\end{tikzpicture}
\end{center}

\begin{center}
\renewcommand{\arraystretch}{1.3}
\begin{tabular}{@{}p{3.1cm}p{11.3cm}@{}}
\toprule
\textbf{Etapa} & \textbf{Qué hace y qué garantía de SOT aporta} \\
\midrule
\term{Source Connector} & Trae bytes crudos de la fuente. El KQL es una ``bomba de extracción'', nunca lógica de hipótesis. Registra cobertura por-fuente (si una tabla falla, se marca ciega). \\
\term{Raw Event Journal} & Log append-only de cada evento tal como llegó, con rotación por tamaño para no llenar el disco. Permite \emph{reconstruir} el grafo desde cero (base de confianza). \\
\term{Normalizer} & Convierte cada evento en un \emph{sobre canónico}: id de fuente, tiempo de evento \emph{y de ingesta}, procedencia, hash, payload crudo. Sin tipos de Azure en el núcleo. \\
\term{GraphDelta} & Contrato de cambios \emph{idempotente}: re-ingestar el mismo evento no duplica aristas. Clave: dedup por \emph{huella canónica}, nunca por timestamp. \\
\term{Incremental Runtime} & Aplica los deltas al grafo en memoria. Garantía verificada: \emph{incremental $=$ rebuild} (aplicar deltas da el mismo grafo que reconstruir desde el journal). \\
\term{Domain API} & Expone \code{run\_hunt} y demás operaciones, agnóstico al transporte. GUI, REST y MCP invocan \emph{el mismo} código local. \\
\term{MCP / GUI} & El consumidor. Recibe los caminos encontrados \textbf{y} los metadatos de cobertura, para nunca confundir ``ciego'' con ``vacío''. \\
\bottomrule
\end{tabular}
\end{center}

% ======================================================================
\section{Paso a paso: la vida de un evento}

Sigamos un único evento --- un login de Windows --- desde que ocurre hasta que
un analista caza sobre él. Este es el recorrido que hace real a SOT.

\begin{enumerate}[leftmargin=*,itemsep=6pt]
  \item[\textbf{Paso 0.}] \term{Ocurre el hecho.} A las 14:30, la usuaria
  \code{alice} se autentica en \code{WORKSTATION-01}. Microsoft Sentinel lo
  registra como un \code{SecurityEvent} con \code{EventID 4624}. Este es el
  estado del mundo real, $X$.

  \item[\textbf{Paso 1.}] \term{Extracción cruda.} El Source Connector polea
  Sentinel con una query \emph{ligera}, sin joins ni agregaciones:
\begin{lstlisting}
SecurityEvent
| extend _ingest_ts = coalesce(ingestion_time(), TimeGenerated)
| where _ingest_ts > datetime(<ultimo watermark>)
| order by _ingest_ts asc | take 10000
\end{lstlisting}
  \begin{warnbox}[Por qué \code{ingestion\_time()} y no \code{TimeGenerated}]
  El tiempo del \emph{evento} (\code{TimeGenerated}) puede llegar ``hacia
  atrás'': un log tardío entra al workspace con un timestamp anterior al
  cursor y se perdería para siempre. El tiempo de \emph{ingesta} solo crece,
  así que el cursor \code{> watermark} nunca se saltea una llegada tardía.
  \end{warnbox}

  \item[\textbf{Paso 2.}] \term{Sobre canónico + journal.} El evento se envuelve
  en un \code{CanonicalEnvelope} que congela: el tiempo de evento
  \emph{y} el de ingesta, la procedencia (qué mapping lo produjo), y un hash
  canónico determinístico. Se escribe al journal append-only.

  \item[\textbf{Paso 3.}] \term{GraphDelta idempotente.} El normalizador emite
  un delta con clave \code{event\_id} $=$ \code{canonical\_hash} $=$
  \code{sha256(payload $\|$ mapping\_hash)}. Si el conector se cae y re-polea
  la misma fila, el delta tiene la \emph{misma} clave $\Rightarrow$ el runtime
  lo reconoce como ya aplicado y \textbf{no duplica la arista}.

  \item[\textbf{Paso 4.}] \term{Aplicación al grafo.} El runtime inserta la
  arista \code{User(alice) --[Auth]--> Host(WORKSTATION-01)}. La arista lleva:
  su tiempo de evento (para ordenar), su tiempo de \emph{ingesta}
  (\code{ingested\_at}, para razonar sobre llegada tardía), y su
  \code{mapping\_hash} \emph{congelado} (para procedencia estable).

  \item[\textbf{Paso 5.}] \term{Fusión multi-fuente.} Otro evento --- una
  conexión de FortiGate desde \code{WORKSTATION-01} a \code{10.0.0.42} ---
  entra por el mismo camino. Como comparte el nodo \code{Host}, las dos fuentes
  se \textbf{fusionan en un solo grafo conectado}. Sentinel es solo una fuente
  más.

  \item[\textbf{Paso 6.}] \term{La caza, 100\% local.} El analista lanza la
  hipótesis \code{User -[Auth]-> Host -[Connect]-> IP}. El \code{run\_hunt}
  \emph{traversa el grafo local} --- sin volver a la nube, sin KQL --- y
  encuentra el camino \code{alice} $\to$ \code{WORKSTATION-01} $\to$
  \code{10.0.0.42}.

  \item[\textbf{Paso 7.}] \term{La respuesta honesta.} La respuesta no es solo
  ``1 camino''. Incluye un objeto \code{AnswerCoverage}:
\begin{lstlisting}
AnswerCoverage {
  possibly_incomplete: false,
  coverage_start: 1705328880, coverage_end: 1705329060,
  missing_sources: []
}
\end{lstlisting}
  Si en cambio el conector de FortiGate hubiera fallado en esa ventana, la
  \emph{misma} caza devolvería el mismo camino \textbf{pero} con
  \code{missing\_sources: ["fortigate:edge-fw"]} y
  \code{possibly\_incomplete: true}. \emph{Ahí} está la diferencia entre
  ``no lo vi'' y ``no pasó''.
\end{enumerate}

% ======================================================================
\section{Los canales: en qué dimensiones puede fallar el observador}

SOT descompone la suficiencia en \term{canales de telemetría}: dimensiones
independientes en las que la observación puede ser fiel o ciega. Cada hipótesis
de caza declara qué canales necesita. GraphHunter tiene 11 canales; siete son
requeridos por las hipótesis del catálogo, y todos están hoy en \textbf{full}.

\begin{center}
\renewcommand{\arraystretch}{1.25}
\begin{tabular}{@{}p{3.6cm}p{10.8cm}@{}}
\toprule
\textbf{Canal} & \textbf{Pregunta que responde} \\
\midrule
\code{entity\_identity} & ¿Dos menciones de ``alice'' son la misma entidad? \\
\code{edge\_multiplicity} & ¿Cuántas veces ocurrió realmente esta relación? (dedup sin perder multiplicidad) \\
\code{event\_time} & ¿Cuándo ocurrió el evento? \\
\code{ingestion\_time} & ¿Cuándo lo \emph{recibimos}? (para razonar sobre llegada tardía) \\
\code{field\_survival} & ¿Sobrevive la fila cruda para reconstruir campos no mapeados? \\
\code{row\_completeness} & ¿Registramos qué filas se descartaron por malformadas? \\
\code{ingest\_completeness} & ¿Sabemos si la ingesta se cortó a mitad (memoria crítica)? \\
\code{source\_coverage} & ¿Sabemos si una fuente quedó ciega en la ventana? \\
\code{skip\_diagnostics} & ¿Guardamos muestras de lo que se saltó? \\
\code{completeness\_flag} & ¿La respuesta puede decir ``posiblemente incompleta''? \\
\code{provenance\_stability} & ¿Editar un mapping después de ingestar contamina la procedencia vieja? \\
\bottomrule
\end{tabular}
\end{center}

\begin{ideabox}[Preguntas transversales]
Además de las 30 hipótesis del catálogo, hay preguntas \emph{transversales} que
combinan canales. La más exigente es \code{loss\_reconstruction}
(``¿podemos reconstruir qué se perdió y por qué?''), que requiere
\code{completeness\_flag}, \code{ingestion\_time}, \code{provenance\_stability}
y \code{skip\_diagnostics}. Con los cuatro en \textbf{full}, esta transversal
es \term{Suficiente}.
\end{ideabox}

% ======================================================================
\section{La auditoría: cómo se \emph{mide} la suficiencia}

La suficiencia no se afirma, se \textbf{audita}. Un runner (que vive en un
repositorio separado, para preservar la cultura de resultados congelados)
evalúa cada hipótesis contra el modelo de canales declarado y emite un
veredicto:

\begin{center}
\begin{tikzpicture}[node distance=4mm,every node/.style={font=\small}]
  \node[draw=ghgreen,fill=ghgreen!8,rounded corners,text width=4.6cm,align=center,minimum height=1.3cm] (suf) {\textbf{Sufficient}\\ conjunto ciego \emph{vacío}: la respuesta es confiable};
  \node[draw=ghred,fill=ghred!8,rounded corners,text width=6.2cm,align=center,minimum height=1.3cm,right=1.4cm of suf] (ins) {\textbf{InsufficientWithWitness}\\ hay un canal ciego $\Rightarrow$ se adjunta el testigo concreto};
\end{tikzpicture}
\end{center}

El veredicto es puramente mecánico: una hipótesis es \emph{Suficiente} si y
solo si su \emph{conjunto ciego} --- los canales requeridos que no están en
\code{full} --- es vacío. Esto hace que cerrar un canal a \code{full}
\textbf{propague} suficiencia de forma verificable y honesta.

\begin{keybox}[Disciplina de claims (tabla de honestidad)]
Cada afirmación de garantía en la documentación lleva una etiqueta:
\emph{proved / tested / implemented / read-verified / proposed}. Un canal se
declara \code{full} solo con evidencia \code{file:line} concreta \emph{y} con
sus límites declarados. Ejemplos de límites que GraphHunter declara sin
maquillar: el reloj de ingesta es el de \emph{recepción} de GraphHunter (no el
del origen); la identidad de un parser builtin es \emph{code-versioned}; la
ventana del journal es reciente-only.
\end{keybox}

% ======================================================================
\section{Dos cierres de canal, en detalle}

Para hacer tangible cómo se cierra un testigo, veamos dos casos reales.

\subsection{\texttt{ingestion\_time}: de ``ciego'' a ``full''}

\textbf{El testigo:} la arista guardaba \code{ingested\_at\_delta = 0}
hardcodeado, así que un evento tardío y uno puntual producían la
\emph{misma} arista --- \code{same\_observed}.

\textbf{El cierre:} (1) los sitios de ingesta ahora estampan
\code{encode\_ingested\_at(now())}; (2) el valor se \emph{preserva} en la
compactación; (3) se \emph{lee de vuelta} en la arista materializada. Ahora un
evento tardío estampa un \code{ingested\_at} estrictamente posterior
$\Rightarrow$ el testigo se rompe. Verificado por un test que ingesta un
evento y verifica que su \code{ingested\_at} $\approx$ ahora, no la época base.

\subsection{\texttt{provenance\_stability}: el testigo que casi se nos escapa}

\textbf{El testigo:} editar un mapping después de ingestar no debe cambiar la
procedencia de aristas viejas. Un análisis inicial concluyó ``ya está cerrado''
--- pero la revisión adversarial encontró un camino de contaminación real: una
función que \emph{sobreescribía} el \code{field\_config} de un dataset ya
ingestado, y el export leía ese valor \emph{mutable}.

\textbf{El cierre:} se congeló la identidad del mapping
(\code{DatasetInfo.mapping\_hash}) \emph{separada} del \code{field\_config}
mutable; el export lee solo el valor congelado. Un test \emph{genuino} ingesta
bajo el config A, edita a config B, re-exporta, y verifica que las aristas
viejas mantienen el hash de A.

\begin{ideabox}[La lección]
Este caso es el espíritu de SOT en acción: la disciplina de \emph{verificar
contra el código, no contra la nota}, y de dejar que una revisión adversarial
cace lo que el autor pasó por alto. La honestidad no es un adorno; es el
mecanismo.
\end{ideabox}

% ======================================================================
\section{El certificado incremental $=$ rebuild}

Hay una grieta silenciosa en toda esta historia. GraphHunter no reconstruye el
grafo desde cero en cada ingesta --- eso sería carísimo. Aplica \emph{deltas}
incrementales a un grafo que ya vive en memoria. Pero entonces surge una duda
que va al corazón de la confianza:

\begin{ideabox}[La pregunta de fondo]
El camino \emph{rápido} (aplicar deltas incrementalmente) ¿produce
\textbf{el mismo grafo} que el camino \emph{lento y confiable} (reconstruir
todo desde el journal)? Si no lo hiciera, toda garantía de cobertura o
suficiencia sobre el grafo vivo sería una ilusión: estaríamos razonando sobre
un artefacto que no es fiel a la evidencia.
\end{ideabox}

\subsection{Validación por traducción: quién es la base de confianza}

GraphHunter adopta aquí un patrón tomado de la verificación de compiladores
(SOT-VM): la \term{validación por traducción}. En vez de \emph{confiar} en que
el runtime incremental es correcto, se lo trata como un ``compilador no
confiable'' y se \textbf{verifica su salida} contra una referencia.

\begin{center}
\begin{tikzpicture}[node distance=4mm,every node/.style={font=\small}]
  \node[draw=ghgreen,fill=ghgreen!8,rounded corners,text width=4.3cm,align=center,minimum height=1.5cm] (ref) {\textbf{Rebuild desde el journal}\\ la \emph{referencia confiable}};
  \node[draw=ghamber,fill=ghamber!10,rounded corners,text width=4.3cm,align=center,minimum height=1.5cm,right=2.2cm of ref] (live) {\textbf{Grafo incremental vivo}\\ el artefacto \emph{no confiable}};
  \node[draw=ghblue,fill=ghlight,rounded corners,text width=5.2cm,align=center,minimum height=1.2cm,below=1.2cm of $(ref)!0.5!(live)$] (cmp) {\textbf{El comparador}\\ \emph{la base de confianza}};
  \draw[-{Latex[length=2mm]},ghgray] (ref.south) -- (cmp.north west);
  \draw[-{Latex[length=2mm]},ghgray] (live.south) -- (cmp.north east);
\end{tikzpicture}
\end{center}

La base de confianza no es el runtime incremental (que puede tener bugs), sino
el \textbf{comparador} más la referencia. Si el comparador dice ``iguales'',
entonces --- para todo lo que el comparador sabe mirar --- el camino rápido es
tan fiel como el lento.

\subsection{Qué compara el comparador (y qué honestamente \emph{no})}

El comparador toma una foto de ambos grafos (\code{to\_snapshot()}, resolviendo
las identidades internas a cadenas de texto para poder compararlas entre
grafos) y contrasta dimensión por dimensión:

\begin{center}
\renewcommand{\arraystretch}{1.2}
\begin{tabular}{@{}p{3.9cm}p{10.5cm}@{}}
\toprule
\textbf{Dimensión} & \textbf{Cómo se compara} \\
\midrule
\code{entities} & Conjunto de ids de entidades. \\
\code{entity\_metadata} & Por id: tipo (por \emph{valor}, no por texto), metadata y \code{dataset\_id}. \\
\code{edges\_multiset} & \emph{Multiset} de aristas \code{(src, rel, dst, ts, meta, dataset)}: la multiplicidad se atrapa por construcción; el orden no se asume. \\
\code{temporal\_window} & El rango \code{[min, max]} de timestamps de aristas. \\
\code{hypothesis\_results} & Una batería de hipótesis corrida en \emph{ambos} grafos; se comparan los caminos como multiset. \\
\bottomrule
\end{tabular}
\end{center}

\begin{keybox}[La honestidad del \texttt{NotComparable}]
Hay dimensiones que el grafo \textbf{no observa}: la procedencia más allá del
\code{dataset\_id}, los scores de anomalía (los produce un pase aparte), el
timestamp MVCC de inserción (no determinístico entre live y rebuild). El
comparador las marca \code{NotComparable} \emph{con su razón} --- nunca como
``iguales''. Afirmar igualdad sobre una dimensión que no se puede observar es
justamente la sobre-afirmación que SOT prohíbe. Del mismo modo, si la batería
de hipótesis no produce ningún camino en ninguno de los dos grafos, el
resultado es \code{NotComparable} (``silencio $\neq$ acuerdo''), no un
\code{Match} vacío.
\end{keybox}

\subsection{El certificado sobre un dataset \emph{real}}

La prueba definitiva no corre sobre datos sintéticos, sino sobre un dataset de
intrusión \textbf{multi-fuente real} (Sysmon $+$ Sentinel $+$ FortiGate),
ingerido por \emph{dos caminos genuinamente distintos}:

\begin{center}
\begin{tikzpicture}[node distance=5mm,every node/.style={font=\footnotesize}]
  \node[draw=ghamber,fill=ghamber!10,rounded corners,text width=5.0cm,align=center,minimum height=1.4cm] (a) {\textbf{Camino \code{live}}\\ en memoria, incremental, en batches por-fuente $+$ un \emph{crash replay}};
  \node[draw=ghgreen,fill=ghgreen!8,rounded corners,text width=5.0cm,align=center,minimum height=1.4cm,right=1.6cm of a] (b) {\textbf{Camino \code{rebuilt}}\\ desde cero, leyendo el journal \emph{de disco}};
  \node[draw=ghblue,fill=ghlight,rounded corners,text width=4.6cm,align=center,below=1.0cm of $(a)!0.5!(b)$,minimum height=1.0cm] (eq) {\code{all\_comparable\_dimensions\_match()}\\ \textcolor{ghgreen}{\textbf{$=$ verdadero}}};
  \draw[-{Latex[length=2mm]},ghgray] (a.south) -- (eq.north west);
  \draw[-{Latex[length=2mm]},ghgray] (b.south) -- (eq.north east);
\end{tikzpicture}
\end{center}

Es una equivalencia real \emph{memoria-vs-disco} $+$ \emph{incremental-vs-batch}
$+$ \emph{dedup} (el crash replay re-aplica un batch y el conteo \textbf{no}
crece). Resultado: 9 entidades / 7 relaciones, las cinco dimensiones
comparables en \code{PASS}, y la procedencia/scores honestamente
\code{NotComparable}.

\begin{warnbox}[El control negativo: el certificado tiene dientes]
Un certificado que siempre dice ``iguales'' no prueba nada. Por eso hay un
\emph{control negativo}: se dropea el último registro del journal, de modo que
la reconstrucción sea un subconjunto estricto del grafo vivo. El comparador
entonces marca \code{DIFF} correctamente. Un certificado que pasa
\emph{y} falla cuando debe fallar es un certificado en el que se puede confiar.
\end{warnbox}

\subsection{Un límite declarado: \texttt{rebuild\_faithful}}

La reconstrucción todavía no reproduce fielmente los sobres de
\code{Correction}/\code{Tombstone} (correcciones y borrados lógicos). Si el
journal validado contiene alguno, el reporte marca
\code{rebuild\_faithful: false} con un caveat --- la equivalencia \textbf{no}
se afirma sobre un journal con correcciones. Otro ejemplo de la disciplina:
declarar el límite en vez de maquillarlo.

\begin{ideabox}[Por qué esto cierra el círculo]
El certificado incremental $=$ rebuild es lo que hace \emph{confiables} a todas
las garantías anteriores. La cobertura, la multiplicidad, la procedencia
congelada --- todas se calculan sobre el grafo vivo (el camino rápido). El
certificado prueba que ese grafo rápido es fiel a la evidencia del journal (el
camino confiable). Sin él, SOT razonaría sobre una sombra; con él, razona sobre
el grafo real.
\end{ideabox}

% ======================================================================
\section{El estado final: todo Suficiente, con límites honestos}

Tras cerrar los canales, el estado de la auditoría es:

\begin{center}
\renewcommand{\arraystretch}{1.3}
\begin{tabular}{@{}p{6.5cm}p{4cm}@{}}
\toprule
\textbf{Familia} & \textbf{Veredicto} \\
\midrule
30 hipótesis del catálogo ATT\&CK & \textcolor{ghgreen}{\textbf{Sufficient}} (30/30) \\
Transversal \code{exhaustiveness\_marking} & \textcolor{ghgreen}{\textbf{Sufficient}} \\
Transversal \code{loss\_reconstruction} & \textcolor{ghgreen}{\textbf{Sufficient}} \\
Los 11 canales de telemetría & \textcolor{ghgreen}{\textbf{full}} \\
\bottomrule
\end{tabular}
\end{center}

Y la maquinaria corre \textbf{por defecto} (journal acotado + dedup idempotente
+ cobertura + cobertura por-fuente), con cada flag como un escape hatch que se
puede apagar. Lo que \emph{no} se afirma: que el repositorio sea ``Suficiente
en todo'' de forma incondicional. Cada garantía arrastra su límite declarado
--- y eso, precisamente, es lo que la hace confiable.

\begin{keybox}[La idea para llevarse]
SOT no le pide al observador que sea omnisciente. Le pide que sea
\textbf{honesto sobre su ceguera}. En GraphHunter, eso se traduce en una regla
simple: \emph{una respuesta de caza nunca presenta datos incompletos como
exhaustivos}. Cuando el sistema no puede ver algo, lo dice --- y esa
distinción, entre ``no pasó'' y ``no lo vi'', es lo que separa a un cazador de
amenazas confiable de uno que solo parece confiable.
\end{keybox}

\vfill
\begin{center}
\small\color{ghgray}\emph{Documento generado a partir de la arquitectura de
GraphHunter y su serie de validación SOT.}
\end{center}

\end{document}
